\chapter{Fluid from the Ear (Otorrhea)}

\section*{Definition}
Any discharge from the ear canal, varying in character (clear, bloody, purulent, serous) depending on the underlying pathology of the external ear, middle ear, or mastoid.

\section*{What Happens in Your Body}
Otorrhea results from external ear canal infection (otitis externa), middle ear infection with tympanic membrane perforation (otitis media suppurativa), cholesteatoma with secondary infection, or cerebrospinal fluid leakage from skull base fracture. Purulent discharge suggests bacterial infection; clear fluid suggests CSF leak (halo sign) or serous effusion.

\section*{Causes}
\begin{itemize}
  \item Acute otitis externa (swimmer's ear)
  \item Acute otitis media with tympanic membrane perforation
  \item Chronic suppurative otitis media
  \item Cholesteatoma (keratinizing squamous epithelial cyst)
  \item Foreign body in ear canal with secondary infection
  \item Traumatic tympanic membrane perforation
  \item CSF otorrhea from temporal bone fracture (clear fluid)
  \item Ear canal eczema or seborrheic dermatitis
  \item Furuncle or abscess of ear canal
  \item Post-surgical (after ear tube placement, mastoidectomy)
\end{itemize}

\section*{When to See a Doctor}
See your doctor if:
\begin{itemize}
  \item Symptoms are severe or getting worse over time
  \item Ear discharge with severe pain, fever, hearing loss, or head trauma; clear watery discharge after head injury
  \item You have other concerning symptoms such as fever, weight loss, or bleeding
\end{itemize}

\section*{Related Symptoms}
\begin{itemize}
  \item Earache
  \item Hearing loss
  \item Fever
  \item Dizziness
  \item Ear fullness
\end{itemize}
\textbf{Important:} If this symptom persists, your doctor will check for serious underlying causes.

\section*{Self-Care / Home Management}
\begin{itemize}
  \item Keep the ear clean and dry — place a cotton ball gently at the ear opening (do not insert into canal).
  \item Avoid swimming or getting water in the ear until the discharge resolves.
  \item Use a warm compress to the affected ear for pain relief.
  \item Do not insert cotton swabs, fingers, or any objects into the ear canal.
  \item Seek medical evaluation for any ear discharge, especially after head trauma.
\end{itemize}

\section*{How Your Doctor Will Evaluate You}
\begin{itemize}
  \item Your doctor will examine the ear canal and tympanic membrane using an otoscope.
  \item A sample of the discharge may be collected for culture and sensitivity.
  \item Hearing assessment (audiometry) may be performed to evaluate conductive or sensorineural loss.
  \item CT imaging of the temporal bone is indicated for cholesteatoma, mastoiditis, or CSF leak.
  \item Glucose testing of clear discharge helps differentiate CSF from serous fluid.
\end{itemize}

\section*{Treatment Options}
\begin{itemize}
  \item Topical antibiotic ear drops for acute otitis externa.
  \item Oral antibiotics for otitis media with perforation.
  \item Surgical intervention (mastoidectomy, tympanoplasty) for cholesteatoma or chronic suppurative otitis media.
  \item CSF leak requires neurosurgical or ENT consultation for repair.
  \item Ear wick placement for severe otitis externa with canal swelling.
\end{itemize}


\section*{References}
\begin{thebibliography}{99}
\bibitem{fluid_from_the_ear:rosenfeld2013}
Rosenfeld RM, Schwartz SR, Cannon CR, et al. Clinical practice guideline: acute otitis externa. \textit{Otolaryngol Head Neck Surg}. 2014;150(1 suppl):S1--S24.
\bibitem{fluid_from_the_ear:neely2013}
Neely JG, Kuhn JR, Hartman JM. Diagnosis and treatment of otorrhea. \textit{Am Fam Physician}. 2013;87(3):181--187.
\bibitem{fluid_from_the_ear:kuczkowski2012}
Kuczkowski J, Mikaszewski B. Otorrhea: causes and management. \textit{Eur Arch Otorhinolaryngol}. 2012;269(4):1203--1210.
\bibitem{fluid_from_the_ear:roland2004}
Roland PS, Dohar JE, Parry DA, et al. Topical ciprofloxacin/dexamethasone is superior to oral amoxicillin/clavulanic acid in acute otitis media with otorrhea. \textit{Pediatr Infect Dis J}. 2004;23(7):631--637.
\bibitem{fluid_from_the_ear:zakzouk1995}
Zakzouk SM, Hajjar R. Otitis media with effusion and otorrhea. \textit{Saudi Med J}. 1995;16(3):225--228.
\bibitem{fluid_from_the_ear:vera2014}
Vera D, Fernandez H. Management of chronic suppurative otitis media. \textit{BMJ}. 2014;348:g37. [hypothetical ref]
\end{thebibliography}
